\unit{Many-electron atoms, configurations and term symbols}
\unitref{Week 8; Assignment 3 Q2; Assignment 4 Q1}
\subsection{Why the hydrogen formula stops working}
Two electrons repel each other. The
many-electron Hamiltonian has the form
\[
\boxed{H=\sum_i\left[\frac{p_i^2}{2m_e}
-\frac{Ze^2}{4\pi\epsilon_0r_i}\right]
+\sum_{i<j}\frac{e^2}{4\pi\epsilon_0r_{ij}}.}
\]
The last term depends on \emph{two}
electron positions, so we cannot solve
each electron independently as if the
other electrons were absent.
An approximate central potential describes
the nucleus screened by the electron cloud.
The one-electron orbitals are labelled
\(n,l\), but their energies are no longer
functions of \(n\) alone.

\subsection{Pauli exclusion and antisymmetric states}
Electrons have spin \(s=1/2\) and
are identical fermions. Interchanging
two electron labels must reverse
the \emph{total} wavefunction:
\[
\Psi(1,2)=-\Psi(2,1).
\]
A two-electron Slater determinant
automatically enforces this:
\[
\Psi_{ab}(1,2)=\frac1{\sqrt2}
[a(1)b(2)-b(1)a(2)].
\]
If \(a=b\), the determinant vanishes:
two electrons cannot occupy the
same spin-orbital.

\subsection{Direct repulsion and exchange}
Evaluate Coulomb repulsion
\(v_{12}=e^2/(4\pi\epsilon_0r_{12})\)
in a Slater determinant. Expanding
the bra and ket yields
\[
\boxed{\langle v_{12}\rangle=J_{ab}-K_{ab}},
\]
where
\[
J_{ab}=\iint\frac{|a(1)|^2|b(2)|^2}{r_{12}}\,
d1\,d2,
\]
\[
K_{ab}=\iint\frac{a^*(1)b(1)b^*(2)a(2)}
{r_{12}}\,d1\,d2
\]
in atomic units. The direct term resembles
ordinary Coulomb repulsion between densities.
The exchange term arises from the
antisymmetric \emph{wavefunction}.
It is not another classical force;
for orthogonal spin labels it vanishes.
Hartree--Fock finds orbitals self-consistently
so that direct and exchange contributions
enter the mean-field equations.

\subsection{How spectroscopic term notation works}
For light atoms with reasonably good
LS (Russell--Saunders) coupling, first
add the individual orbital momenta to
\(L\), individual electron spins to
\(S\), then couple \(L,S\) to total
electronic \(J\). Write
\[
\boxed{{}^{2S+1}L_J},\qquad
L=0,1,2,3,4\leftrightarrow S,P,D,F,G.
\]
For example, \({}^3F_2\) means
total spin \(S=1\) (triplet),
total orbital angular momentum \(L=3\)
and total electronic \(J=2\).
The allowed \(J\) values are
\(J=|L-S|,|L-S|+1,\ldots,L+S\).

\subsection{Equivalent electrons require extra care}
Two \(d\) electrons each have \(l=2\)
and two spin projections. There are
ten distinct spin-orbitals, hence
\[
\binom{10}{2}=45
\]
antisymmetric configurations.
Without Pauli restrictions, angular-momentum
addition would suggest \(L=0,1,2,3,4\),
\(S=0,1\). But the orbital-coupled
wavefunction has exchange phase
\((-1)^{2l-L}=(-1)^L\).
The spin singlet \(S=0\) is antisymmetric,
while triplet \(S=1\) is symmetric.
Thus only \emph{even \(L\), \(S=0\)}
or \emph{odd \(L\), \(S=1\)} survive:
\[
\boxed{{}^1S_0,\ {}^1D_2,\ {}^1G_4,\quad
{}^3P_{0,1,2},\ {}^3F_{2,3,4}.}
\]

\begin{worked}{Check that every \(d^2\) state is counted}
A term with total \(J\) has \(2J+1\)
magnetic substates. Sum:
\[
1+5+9+(1+3+5)+(5+7+9)=\boxed{45}.
\]
The count agrees with the Pauli
spin-orbital count \(\binom{10}{2}\).
Without this check we could easily
include forbidden LS terms or miss
part of an allowed multiplet.
\end{worked}

\subsection{Hund's rules and why they are only trends}
Within their usual domain, Hund's rules
favour (1) largest total \(S\), then
(2) largest \(L\) for that multiplicity,
then (3) smallest \(J\) in a
less-than-half-filled subshell.
For \(d^2\), this predicts
the ground level \({}^3F_2\).
However, these rules do \emph{not}
unambiguously order every excited term
of a real multi-electron atom.
Titanium's neutral \(3d^24s^2\)
configuration demonstrates the broad
agreement and exceptions when compared
with NIST experimental level data.

\begin{checkpoint}
The notation \({}^{2S+1}L_J\) is a
coupling scheme, not a claim that
electrons orbit independently.
First impose antisymmetry, then
enumerate LS terms, attach \(J\),
and check the total degeneracy.
\end{checkpoint}
